\section{Le giunzioni cellulari}

% ---------------------------------------------------------------------------
\subsection{La matrice extracellulare}

\textbf{Matrice extracellulare} (ECM): rete organizzata di materiale extracellulare presente nelle immediate vicinanze della membrana plasmatica.

\rubrica{Composizione}
\begin{itemize}
  \item \textbf{collagene}: proteina fibrosa;
  \item \textbf{fibronectine}: glicoproteine della ECM che si legano alle integrine e ad altri recettori presenti nella membrana plasmatica;
  \item \textbf{laminina};
  \item \textbf{proteoglicano}: enormi complessi proteina--polisaccaride che occupano la maggior parte del volume dello spazio extracellulare, con siti di legame reciproco e per recettori (integrine) localizzati sulla superficie delle cellule;
  \item \textbf{integrina}: recettore di membrana che lega la ECM al citoscheletro (filamento intermedio, microfilamenti).
\end{itemize}

\rubrica{Caratteristiche}
\begin{itemize}
  \item le proteine presenti nella matrice extracellulare servono come impalcature, travi, cemento e reti;
  \item può assumere diverse forme nei differenti tessuti;
  \item la maggior parte delle proteine presenti sono di tipo fibroso;
  \item può avere un ruolo regolatorio nel determinare la forma e le attività di una cellula.
\end{itemize}

% ---------------------------------------------------------------------------
\subsection{Lamina basale}

Funzioni:
\begin{itemize}
  \item fornisce un supporto meccanico per l'adesione cellulare;
  \item genera segnali per la sopravvivenza cellulare;
  \item serve da substrato per la migrazione cellulare;
  \item separa tessuti adiacenti all'interno di un organo;
  \item funziona da barriera al passaggio di macromolecole.
\end{itemize}

% ---------------------------------------------------------------------------
\subsection{Panoramica delle giunzioni cellulari}

Quattro tipi principali: \textbf{giunzione gap} (comunicante), \textbf{giunzioni aderenti}, \textbf{giunzioni strette}, \textbf{desmosoma}.
\begin{itemize}
  \item \textbf{giunzione gap}: connessoni;
  \item \textbf{giunzioni aderenti}: caderina, filamenti di actina;
  \item \textbf{giunzioni strette}: complessi proteici;
  \item \textbf{desmosoma}: cadherin, filamenti di actina.
\end{itemize}

% ---------------------------------------------------------------------------
\subsection{Desmosomi}

Ciascuna delle due cellule compartecipa alla formazione del desmosoma. Ciascun desmosoma è costituito da un paio di dischi a forma di bottone associati con le membrane plasmatiche di cellule adiacenti e da filamenti proteici intercellulari che connettono queste cellule. I filamenti intermedi attaccati a questi dischi sono connessi con altri desmosomi.

Poiché sono ancorati su entrambi i lati cellulari, risultano essere molto robusti.

\rubrica{Struttura}
\begin{itemize}
  \item struttura densa (\textbf{placca}) sul versante citoplasmatico della membrana;
  \item molecole di adesione (\textbf{desmogleina}, \textbf{desmocollina}) si estendono attraverso lo spazio intercellulare fino all'altra cellula, in corrispondenza delle proteine della placca;
  \item a livello citoplasmatico sono presenti proteine del citoscheletro, in particolare la cheratina;
  \item nella placca: \textbf{placoglobina} e \textbf{desmoplachina}, a cui si ancora il filamento intermedio.
\end{itemize}

% ---------------------------------------------------------------------------
\subsection{Emidesmosomi}

Siti differenziati alla superficie basale delle cellule epiteliali, dove le cellule sono attaccate alla membrana basale sottostante.

\begin{itemize}
  \item le molecole di \textbf{integrina} $\alpha_6\beta_4$ delle cellule epidermiche sono collegate ai filamenti intermedi citoplasmatici da una proteina chiamata \textbf{plectina}, presente nelle placche dense, e alla membrana basale da filamenti di ancoraggio formati da un tipo particolare di laminina (\textit{laminina-5});
  \item negli emidesmosomi è presente anche una seconda proteina transmembrana (\textbf{BP180});
  \item le fibre di collagene (collagene tipo VII) fanno parte del derma sottostante.
\end{itemize}

% ---------------------------------------------------------------------------
\subsection{Le giunzioni strette (tight junctions)}

\begin{itemize}
  \item le membrane plasmatiche adiacenti giungono in intimo contatto tra loro, annullando lo spazio intercellulare;
  \item lungo la zona di contatto gli strati esterni delle 2 membrane si fondono, formando un'unica struttura \textbf{pentalaminare} data dalla fusione di proteine intrinseche di membrana;
  \item una giunzione serrata si forma per mezzo di connessioni tra file di proteine di membrana adiacenti; tali proteine sono strettamente impacchettate in file che sigillano lo spazio intercellulare, impedendo il passaggio di materiali negli spazi tra le cellule.
\end{itemize}

Impediscono il passaggio di materiali attraverso gli spazi intercellulari.

% ---------------------------------------------------------------------------
\subsection{Le giunzioni comunicanti (gap junctions)}

Permettono il trasferimento di piccole molecole e ioni tra cellule adiacenti.

\begin{itemize}
  \item connessioni di tipo proteico (\textbf{connessoni}) legano le due membrane plasmatiche;
  \item le particelle proteiche sono attraversate da un canale idrofilo che permette lo scambio di ioni e piccole molecole tra le due cellule;
  \item le due membrane plasmatiche contengono cilindri costituiti da sei molecole di \textbf{connessina}; i due cilindri di membrane opposte sono uniti a formare un canale che connette i compartimenti citoplasmatici delle due cellule;
  \item il connessone (cilindro) è formato da subunità di connessina;
  \item il poro di una giunzione comunicante può aprirsi e chiudersi (conformazioni \textit{chiuso}/\textit{aperto});
  \item le proteine sono distribuite a distanze regolari, formando tra le parti sporgenti una rete di spazi che separa le due membrane adiacenti ed è accessibile dagli interstizi intercellulari circostanti.
\end{itemize}

% ---------------------------------------------------------------------------
\subsection{Le giunzioni aderenti}

L'actina e le caderine permettono alle membrane di essere strettamente legate in modo da impedirne la disgiunzione.

\begin{itemize}
  \item \textbf{caderina}: proteina transmembrana di adesione, extracellulare;
  \item \textbf{$\alpha$-actinina}: lega il filamento di actina alla $\beta$-catenina;
  \item la caderina si lega, sul versante citoplasmatico, a \textbf{$\beta$-catenina}, \textbf{$\alpha$-catenina} e alla \textbf{catenina p120};
  \item la $\beta$-catenina è collegata a una proteina che lega l'actina, ancorando il complesso al filamento di actina;
  \item i segnali associati al complesso di giunzione sono trasmessi al citoplasma e al nucleo.
\end{itemize}

% ---------------------------------------------------------------------------
\subsection{I plasmodesmi}

Mettono in comunicazione il citoplasma di cellule adiacenti, permettendo il passaggio di acqua, ioni e piccole molecole.

\begin{itemize}
  \item internamente i canali sono rivestiti dalle membrane plasmatiche fuse delle due cellule adiacenti;
  \item il \textbf{desmotubulo} consiste in una membrana in continuità con il reticolo endoplasmatico (RE) di entrambi i lati della membrana plasmatica;
  \item l'\textbf{annulus} è la regione anulare del plasmodesma attorno al desmotubulo;
  \item passano attraverso l'\textit{annulus} le molecole che si spostano da una cellula all'altra;
  \item il trasporto di proteine e mRNA attraverso i plasmodesmi è coinvolto nella localizzazione tissutale di alcuni fattori di trascrizione.
\end{itemize}
